\paragraph{} For all NFPA 70 requirements the 2014 draft version was used.

\begin{itemize}
	\item Lockout/Tagout Procedure - Annex G {\newline} {\indent}
	      We followed the sequence of the lockout/tagout procedure as relevant.
	\item Conductor sizing - 240.4(D)(1) - 310.15 {\newline} {\indent}
	      18 gauge wire was used in the panel according with table 310.15(B)(16).
	      This decision was based on an  expected expected maximum current draw of 0.982A for this control system.
	      And is well under the maximum 7A that 18 gauge wire is allowed.
	      % 0.9(power supply)+16m(PLC)+16m(inputs)+10m(outputs)+40m(both lights)=982mA	
	\item Ampacity factor - 210.19(A)(1) - 210.20(A) {\newline} {\indent}
	      Because this system will be operated as a continuous load the ampacity of the conductors used is not less
	      than the noncontinuous load plus 125 percent of the continuous load. Making the minimum allowable ampacity
	      1.22A, well within the 7A range.
	      % 0.9(power supply)+16m(PLC)+16m(inputs)+10m(outputs)+40m(both lights)=982mA	
	\item Fuse sizing - 240.6 {\newline} {\indent}
	      Based on the max expected current draw of 0.982A and the NFPA paragraph listed we selected 1A fuses.
	\item Equipment grounding - 250.119 {\newline} {\indent}
	      All grounding wires are green in compliance with the paragraph listed.
	\item Conductor location use - 310.10(A) {\newline} {\indent}
	      All wires selected for this system are of type THHN. This wire is acceptable for the dry conditions in which
	      our panel will operate.
\end{itemize}
